\section{Posa't a prova amb un test general senzill!}
\iniciTest{t3}{d,a,b,d,a,a,b,b,b,c}

\begin{enumerate}
\pregunta Considera l'enunciat: "Si un nombre és parell ($p$), llavors el seu
quadrat és parell ($q$)." Quina de les següents afirmacions és falsa?

\begin{enumerate}
\opcio{a}{ La seva contrarecíproca és: "Si el quadrat d'un nombre no
és parell, llavors el nombre no és parell."

}
\opcio{b}{ Si $p$ és cert, $q$ també ha de ser cert.

}
\opcio{c}{ Si $q$ és fals, $p$ també ha de ser fals.

}
\opcio{d}{ L'enunciat és vertader, però la seva recíproca "Si el
quadrat d'un nombre és parell, el nombre és parell" és falsa.
}
\end{enumerate}
\feedback

\pregunta La proposició $(p\vee q)\rightarrow(\lnot p\wedge r)$ és falsa quan:

\begin{enumerate}
\opcio{a}{ $p$ és vertadera, $q$ és falsa, i $r$ és falsa.

}
\opcio{b}{ $p$ és falsa, $q$ és vertadera, i $r$ és vertadera.

}
\opcio{c}{ $p$ és falsa, $q$ és falsa, i $r$ és falsa.

}
\opcio{d}{ $p$ i $q$ són falses, i $r$ és vertadera.
}
\end{enumerate}
\feedback

\pregunta Suposem que l'univers és el conjunt dels nombres enters ($\mathbb{Z}$).
Si $P(x)$ és "$x$ és un nombre parell" i $Q(x)$ és "$x$ és un
nombre senar", quina de les següents afirmacions és vertadera?

\begin{enumerate}
\opcio{a}{ $\forall x(P(x)\wedge Q(x))$ és vertadera.

}
\opcio{b}{ $\exists x(P(x)\wedge Q(x))$ és falsa.

}
\opcio{c}{ $\lnot\exists xP(x)$ és equivalent a "Cap nombre enter és senar."

}
\opcio{d}{ $\forall x(P(x)\vee Q(x))$ és falsa.
}
\end{enumerate}
\feedback

\pregunta Si l'afirmació "Si plou ($p$), llavors el carrer es mulla ($q$)"
és falsa, quina de les següents conclusions és certa?

\begin{enumerate}
\opcio{a}{ Està plovent i el carrer està mullat.

}
\opcio{b}{ No plou i el carrer no està mullat.

}
\opcio{c}{ No plou, però el carrer està mullat.

}
\opcio{d}{ Està plovent, però el carrer no està mullat.
}
\end{enumerate}
\feedback

\pregunta Per demostrar que si $n^{2}$ és senar, llavors $n$ és senar (per
a qualsevol enter $n$), la demostració per contrarecíproc ha de
provar que:

\begin{enumerate}
\opcio{a}{ Si $n$ és parell, llavors $n^{2}$ és parell.

}
\opcio{b}{ Si $n$ és senar, llavors $n^{2}$ és senar.

}
\opcio{c}{ Si $n^{2}$ és parell, llavors $n$ és parell.

}
\opcio{d}{ Si $n^{2}$ és senar, llavors $n$ és senar.
}
\end{enumerate}
\feedback

\pregunta Per provar que la suma de dos nombres senars consecutius és sempre
un múltiple de 4, el primer pas de la demostració directa seria
expressar els dos nombres com:

\begin{enumerate}
\opcio{a}{ $2k+1$ i $2k+3$ per a algun enter $k$.

}
\opcio{b}{ $2k$ i $2k+2$ per a algun enter $k$.

}
\opcio{c}{ $n$ i $n+1$ per a algun enter $n$.

}
\opcio{d}{ $k$ i $k+2$ per a algun enter $k$.
}
\end{enumerate}
\feedback

\pregunta La demostració de què $\sqrt{3}$ és irracional utilitzant la
reducció a l'absurd comença suposant que:

\begin{enumerate}
\opcio{a}{ $\sqrt{3}$ és un nombre enter.

}
\opcio{b}{ $\sqrt{3{}}$ és un nombre racional.

}
\opcio{c}{ $\sqrt{3{}}=p/q$, on $p$ i $q$ són parells.

}
\opcio{d}{ $\sqrt{3{}}=p/q$, on $p$ i $q$ tenen factors comuns.
}
\end{enumerate}
\feedback

\pregunta Quina de les següents parelles de proposicions són equivalents?

\begin{enumerate}
\opcio{a}{ $p\vee q$ i $\lnot(\lnot p\vee\lnot q)$.

}
\opcio{b}{ $p\rightarrow q$ i $\lnot p\vee q$.

}
\opcio{c}{ $p\wedge q$ i $\lnot(\lnot p\wedge\lnot q)$.

}
\opcio{d}{ $p\leftrightarrow q$ i $(p\rightarrow q)\vee(q\rightarrow p)$.
}
\end{enumerate}
\feedback

\pregunta Identifica la fal·làcia en el següent
raonament: "Si un estudiant estudia ($p$), aprova l'examen ($q$). L'estudiant
ha aprovat l'examen ($q$). Per tant, l'estudiant ha estudiat ($p$)."

\begin{enumerate}
\opcio{a}{ Fal·làcia de negació de l'antecedent.

}
\opcio{b}{ Fal·làcia d'afirmació del conseqüent.

}
\opcio{c}{ Reducció a l'absurd.

}
\opcio{d}{ Sil·logisme disjuntiu.
}
\end{enumerate}
\feedback

\pregunta L'enunciat "No hi ha cap estudiant que sigui bo en matemàtiques i en
física" es pot expressar formalment com:

\begin{enumerate}
\opcio{a}{ $\exists x(M(x)\wedge F(x))$

}
\opcio{b}{ $\forall x(\lnot M(x)\wedge F(x))$

}
\opcio{c}{ $\lnot\exists x(M(x)\wedge F(x))$

}
\opcio{d}{ $\forall x(\lnot M(x)\wedge\lnot F(x))$
}
\end{enumerate}
\feedback
\end{enumerate}